\section{两阶段推理与键值缓存}
\label{l12:sec:cache}

\subsection{预填充与增量解码的时序}

处理一个新的文本输入时，初始的多个词元都已知。预填充（prefill）把整个初始前缀转换为输入矩阵，并一次经过全部网络层。各位置虽然只能访问因果掩码允许的前缀，但可以一起进行投影、矩阵乘法和逐行归一化。最后一个位置的输出分布用于选择第一个新词元。

后续的增量解码（incremental decoding）每次只把刚生成的一个词元作为新的输入位置送入网络，同时复用历史位置已计算的状态。设初始前缀为 $t_0,\ldots,t_{N-1}$。预填充完成后，每层都已经得到这 $N$ 个位置的键和值；根据位置 $N-1$ 的最终输出，选择词元 $t_N$。此时词元 $t_N$ 刚被选出，还没有经过网络，它的各层键和值尚未形成。

下一次前向传播把 $t_N$ 放在绝对位置 $N$，逐层计算该新位置的表示和键值，并预测 $t_{N+1}$。所以“得到词元编号”与“形成该词元的各层缓存”发生在相邻的两个计算阶段。词元编号已经由词表选择直接得到，可以马上用于嵌入查表；生成回路不需要先把编号转为文本再分词。

\begin{table}[htbp]
\centering
\caption{初始前缀含 $N$ 个词元时的前三次前向传播。缓存长度只统计已经输入网络的词元。}
\label{l12:tab:decode-timeline}
\begin{tabularx}{\linewidth}{@{}L{2.2cm}Y L{3.0cm}L{2.6cm}@{}}
\toprule
阶段 & 本次送入网络的词元 & 缓存行数 & 本次新词元\\
\midrule
预填充 & $t_0,\ldots,t_{N-1}$ & $N$ & $t_N$\\
增量第1步 & $t_N$，位置为 $N$ & $N+1$ & $t_{N+1}$\\
增量第2步 & $t_{N+1}$，位置为 $N+1$ & $N+2$ & $t_{N+2}$\\
\bottomrule
\end{tabularx}
\end{table}

预填充对给定上下文进行表示计算，这里的阶段名称不表示模型必须执行文本摘要任务。聊天、续写、翻译等生成任务都可以采用同样的预填充与解码流程。

\subsection{历史状态为什么可以复用}

键值缓存（key-value cache，KV cache）保存已经计算出的历史键和值。它的正确性来自因果依赖，以及层内其他运算的逐位置性质。在模型参数固定、位置编号保持不变的确定性推理中，把未来词元追加到序列末尾，不会改变已有位置的表示。

这个结论可以按网络层数证明。首先，位置 $i$ 的初始向量由 $t_i$ 和位置 $i$ 的嵌入相加得到。扩展序列不改变这两个量，所以旧位置的 $X^{(0)}_{i,:}$ 保持不变。假设第 $\ell$ 层所有旧位置的输入已经保持不变，考察该层位置 $i$ 的计算。

层归一化只读取位置 $i$ 的特征，因此它的结果不变，随后得到的 $q_i,k_i,v_i$ 也不变。因果注意力仅使用位置 $0,\ldots,i$ 的键和值；新增的未来位置被屏蔽，因此第 $i$ 行有效分数、归一化分母和加权输出都不变。输出投影、残差相加和前馈网络继续独立处理这一行，使第 $\ell+1$ 层的旧位置输入仍然不变。由此逐层归纳，所有层中历史位置的键和值都可以保存后复用。

缓存必须按层保存，因为第 $\ell$ 层的键和值由该层自己的输入与投影权重产生。只保存初始嵌入无法代替中间层的键值，只保存最后一层也不能支持前面各层的注意力。标准多头结构中，每层保存 $K^{(\ell)},V^{(\ell)}\in\mathbb{R}^{B\times h\times N\times d}$；不同层的数值不同，尽管形状相同。

未来位置使用的是新的查询，它需要与历史键进行匹配，再对历史值做加权和。旧查询只用于生成旧位置的输出，因此标准增量解码不需要缓存历史 $Q$。同样，也不需要保存全部历史注意力概率矩阵。缓存 K/V 保留了后续新查询所需的历史信息，同时避免存储会随位置数平方增长的概率状态。

如果修改已有前缀、改变已有位置编号或替换模型参数，上述不变性条件就发生变化，受影响的缓存需要重新建立。这一条件也说明：缓存的有效性取决于生成时的实际计算依赖。

\subsection{一个新增位置在各层中的计算}

设某层、某头已经缓存 $N$ 个位置。对于新位置 $N$，令 $x_N\in\mathbb{R}^{1\times D}$ 表示送入该头投影的归一化输入。只对这一行做三个投影，并将新键和值追加到对应的历史缓存：
\begin{equation}
\begin{aligned}
 q_N&=x_NW_Q,\qquad k_N=x_NW_K,\qquad v_N=x_NW_V,\\
 \widehat K&=\begin{bmatrix}K_{0:N,:}\\k_N\end{bmatrix},\qquad
 \widehat V=\begin{bmatrix}V_{0:N,:}\\v_N\end{bmatrix},\\
 o_N&=\operatorname{softmax}\!\left(\frac{q_N\widehat K^{\mathsf T}}{\sqrt d}\right)\widehat V.
\end{aligned}
\label{l12:eq:cached-attention}
\end{equation}
切片 $0:N$ 表示包含 $0$、不包含 $N$，即原来的 $N$ 行。新查询的形状为 $1\times d$，扩展后的缓存为 $(N+1)\times d$，分数和权重为 $1\times(N+1)$，输出为 $1\times d$。当前新位置自己的键和值已经包含在缓存中，因此对角线也被正确保留。

当一次只新增一个位置，且缓存仅含当前与过去位置时，缓存中的全部列都满足因果条件，不再需要显式构造一个三角掩码。若存在填充位置或其他可见性限制，仍需按照这些限制屏蔽相应列。若一次追加多个位置，则每个新增位置的可见列数不同，仍需要带历史长度偏移的因果掩码。

第~\ref{l12:sec:cuda} 节的核函数使用方阵中的局部行号作为因果边界。单词元解码的分数矩阵只有一行，但这一行对应的绝对位置是 $N$，其可见范围是 $0,\ldots,N$。因此，把这个长向量作为局部第零行直接套用方阵核函数，只会错误地保留第一列。解码核函数需要使用实际可见长度，或者显式加入历史长度偏移。

所有头只计算新行之后，拼接它们的输出，再经过 $W_O$、残差、归一化和前馈网络，得到下一层的新行。下一层使用自己保存的 K/V 重复相同过程。完成全部层后，最后归一化与词表投影只需处理这个新位置，从而得到再下一个词元的分布。

第~\ref{l12:subsec:example} 节可以直接展示这种复用。先处理前两个位置，缓存键 $(1,0),(0,1)$ 和值 $(1,0),(0,2)$。追加第三个位置时，只需新的 $q_2=(\sqrt2,\sqrt2)$、$k_2=(1,1)$、$v_2=(3,1)$。新查询与三行键的缩放分数为 $(1,1,2)$，因此新输出为 $(1.940292,1)$，与一次处理全部三行的结果相同。前两个输出仍然是 $(1,0)$ 和 $(0.268941,1.462117)$，无需重新计算。

\begin{figure}[htbp]
\centering
\begin{tikzpicture}[
  font=\small,>=Latex,
  box/.style={draw=UIUCBlue,line width=0.7pt,rounded corners=1pt,minimum height=0.72cm,align=center},
  flow/.style={->,line width=0.75pt,draw=UIUCBlue}
]
\node[box,minimum width=2.4cm] (q) at (0,1.8) {新查询 $q_N$\\$1\times d$};
\draw[draw=UIUCBlue,line width=0.7pt,fill=white] (2.5,0.85) rectangle (5.5,3.0);
\draw[draw=UIUCBlue,line width=0.7pt,fill=UIUCOrange!12] (2.5,0.85) rectangle (5.5,1.55);
\node[align=center] at (4,2.35) {历史键 $K_{0:N,:}$\\$N\times d$};
\node at (4,1.2) {新键 $k_N$};
\draw[draw=UIUCBlue,line width=0.7pt,fill=white] (6.9,0.85) rectangle (9.9,3.0);
\draw[draw=UIUCBlue,line width=0.7pt,fill=UIUCOrange!12] (6.9,0.85) rectangle (9.9,1.55);
\node[align=center] at (8.4,2.35) {历史值 $V_{0:N,:}$\\$N\times d$};
\node at (8.4,1.2) {新值 $v_N$};
\node[box,minimum width=3.2cm] (score) at (4,-0.4) {$q_N\widehat K^{\mathsf T}/\sqrt d$};
\node[box,minimum width=3.2cm] (prob) at (4,-1.7) {归一化权重};
\node[box,minimum width=3.0cm] (agg) at (8.4,-1.7) {加权求和};
\node[align=center] (out) at (8.4,-2.9) {新输出 $o_N$\\$1\times d$};
\draw[flow] (q.south)|-(score.west);
\draw[flow] (4,0.85)--(score.north);
\draw[flow] (score)--node[right] {$1\times(N+1)$}(prob);
\draw[flow] (prob)--(agg);
\draw[flow] (8.4,0.85)--node[right] {$\widehat V$}(agg.north);
\draw[flow] (agg)--(out);
\end{tikzpicture}

\caption{单层单头的增量注意力。历史键值保留，新键值追加为最后一行；只用新查询计算一个长度为 $N+1$ 的权重向量，再形成新输出。}
\label{l12:fig:cache-flow}
\end{figure}

\subsection{节省的计算与增加的状态容量}

比较两种实现处理到相同总长度 $N$ 时，产生下一词元分布所做的工作。直接重算方式把全部 $N$ 个位置重新输入网络；缓存方式只处理新增的一行，并访问已有键值。单头的主要差别如表~\ref{l12:tab:cache-cost} 所示。

\begin{table}[htbp]
\centering
\caption{单头注意力在总长度 $N$ 时的计算与持久状态。省略不影响阶数的常数。}
\label{l12:tab:cache-cost}
\begin{tabularx}{\linewidth}{@{}Y L{4.2cm}L{4.2cm}@{}}
\toprule
对象 & 重算完整前缀 & 增量解码\\
\midrule
查询、键和值投影 & $\Theta(NDd)$ & $\Theta(Dd)$\\
分数与值聚合 & $\Theta(N^2d)$ & $\Theta(Nd)$\\
注意力权重的计算行数 & 全部 $N$ 行 & 新增的1行\\
持久注意力状态 & 不保存历史状态 & $2Nd$ 个键值元素\\
\bottomrule
\end{tabularx}
\end{table}

缓存把重复计算历史位置的工作去掉了，但新查询仍要访问全部可见键值。因此，每步的注意力计算和缓存读取量仍随 $N$ 线性增长。所有头的 $hd=D$，所以每层的新键值投影为 $\Theta(D^2)$ 量级，注意力匹配与聚合为 $\Theta(ND)$ 量级。输出投影与前馈网络也只计算新行，不再对历史行重复执行。

缓存容量由两个矩阵、批量、层数、上下文长度和特征宽度共同决定。在标准多头注意力中，若每个缓存元素占 $b_{\mathrm{elem}}$ 字节，则
\begin{equation}
  n_{\mathrm{KV}}=2BLhNd=2BLND,\qquad
  M_{\mathrm{KV}}=2BLND\,b_{\mathrm{elem}}\ \text{字节}.
  \label{l12:eq:kv-memory}
\end{equation}
这个公式假定每个查询头都有对应的独立键和值头。它只统计已经使用的缓存元素；按最大长度预分配的实现，应把 $N$ 替换为分配容量对应的长度。模型权重、其他激活和临时工作区还会额外占用显存。

例如，取 $B=1$、$L=32$、$D=4096$、$h=32$、$d=128$、$N=4096$，并用 $16$ 位数保存缓存，即每个元素占 $2$ 字节，则
\begin{equation}
  M_{\mathrm{KV}}
  =2\times1\times32\times4096\times4096\times2
  =2\,147\,483\,648\ \text{字节}
  =2\ \mathrm{GiB}.
  \label{l12:eq:kv-example}
\end{equation}
这里 $1\ \mathrm{GiB}=2^{30}$ 字节。这是一组说明容量关系的标准多头配置。保持其他参数不变时，上下文长度或并行序列数加倍，已使用的缓存容量也加倍；网络层数增加同样线性增加缓存。

工程实现可以为每层预分配足够的键值空间，并用当前有效长度指示已写入部分。追加时只写新行，避免每步重新分配并复制全部历史缓存。若底层采用分块或分页存储，其逻辑内容仍然是式~\eqref{l12:eq:cached-attention} 中按位置排列的历史键值。


\subsection{变长序列与键值头共享}

一个批量中的不同序列通常具有不同的有效长度。设第 $b$ 个序列的长度为 $N_b$，则标准多头结构实际保存的键值数据为
\begin{equation}
 M_{\mathrm{KV}}=2Lhd\,b_{\mathrm{elem}}\sum_{b=1}^{B}N_b.
 \label{l12:eq:variable-kv}
\end{equation}
如果为每个序列按统一上限 $N_{\mathrm{cap}}$ 预分配，则分配量为 $2BLhN_{\mathrm{cap}}d\,b_{\mathrm{elem}}$，可能大于有效数据量。按块分配可以减少未使用空间；块表、对齐以及最后一块中的空余位置仍会产生额外开销。不同请求的历史词元和中间表示通常不同，因此不能把一个请求的缓存直接用作另一个请求的缓存。

标准多头注意力为每个查询头配置独立的键值头。也可以让多个查询头共用一组键和值。设查询头数为 $h_q$、键值头数为 $h_{\mathrm{kv}}$，每个头的维度为 $d$，模型宽度仍为 $D=h_qd$。等长批量的缓存容量统一写为
\begin{equation}
 M_{\mathrm{KV}}=2BLNh_{\mathrm{kv}}d\,b_{\mathrm{elem}}.
 \label{l12:eq:grouped-kv}
\end{equation}
多头注意力（MHA）取 $h_{\mathrm{kv}}=h_q$；分组查询注意力（grouped-query attention，GQA）取 $1<h_{\mathrm{kv}}<h_q$，通常让 $h_{\mathrm{kv}}$ 整除 $h_q$，每组查询头共享一个键值头；多查询注意力（multi-query attention，MQA）取 $h_{\mathrm{kv}}=1$，所有查询头共享键和值。各查询头仍保留自己的查询，计算各自的匹配分数和输出。

在式~\eqref{l12:eq:kv-example} 的配置中，$32$ 个查询头若对应 $8$ 个键值头，缓存由 $2\ \mathrm{GiB}$ 降至 $0.5\ \mathrm{GiB}$。同一组内的查询还可以复用读入的键值数据，实际复用效果取决于核函数组织。这些结构使用不同形状的投影参数，必须与模型及预训练权重匹配；既定 GPT-2 权重的推理仍采用其原有多头结构。
